\NeedsTeXFormat{LaTeX2e}
\documentclass[a4paper,fontsize=12pt,listof=totoc,bibliography=totoc]{scrbook}

% Add the additional packages and macros that you want to use
% Here you will find some useful as well as required packages that are automatically imported in this template with a brief description what they are used for 
%---------------------------------
% Minimally required packages 
%---------------------------------
% Required for correct character encoding
\usepackage[utf8]{inputenc}

\setkomafont{disposition}{\normalcolor\bfseries}
\usepackage[british]{babel}
% To customize page layout
\usepackage{graphicx}
\usepackage{scrhack}
\usepackage{geometry}
\usepackage{layout}
\usepackage{changepage}
\usepackage[final]{microtype}
\usepackage{tikz}
\usetikzlibrary{calc}
% To represent URLs within LaTex
\usepackage{url}
% Check this package
%\usepackage[multiuser,marginclue,nomargin,inline,index]{fixme}
% Needed to include images within LaTex
\usepackage{graphicx}

% Make sure that the numbering of figures and tables is continuous and not dependent on the chapter
\counterwithout{figure}{chapter}
\counterwithout{table}{chapter}
\counterwithout{equation}{chapter}


% Required to represent code within LaTex
\usepackage{algpseudocode}
\usepackage{algorithm}
\usepackage{listings}
\usepackage{xcolor}

\lstdefinestyle{mystyle}{
    basicstyle=\ttfamily\footnotesize,
    breakatwhitespace=false,         
    breaklines=true,                 
    captionpos=b,                    
    keepspaces=true,                 
    numbers=left,                    
    numbersep=5pt,                  
    showspaces=false,                
    showstringspaces=false,
    showtabs=false,                  
    tabsize=2
}

\lstset{style=mystyle}

% Required for bibliography style and inclusion of the bib file
\usepackage{natbib}
\usepackage[a-2u]{pdfx}
%\usepackage{hyperref}
%---------------------------------


%---------------------------------
% Additional packages that can be changed
%%---------------------------------
% Facilitates using colors within LaTex
\usepackage{xcolor}
% To create and work with glossaries
\usepackage[acronym]{glossaries}
%---------------------------------

\usepackage{multirow}

%For HCW master's thesis, please include the following title page and declaration of authorship
% -- Commands for the title page --------------

% -- Title of the work
\newcommand{\TitleValue}{Title}
\newcommand{\SubtitleValue}{Subtitle}

% -- Name
\newcommand{\NameValue}{\textgreater degree(s) first name family name degree(s)\textless \par }


% -- Desired degree
\newcommand{\DegreeValue}{Master of Science (MSc \hspace{0.1cm}(MSC)}

% -- Year
\newcommand{\YearValue}{\textgreater Year\textless }

% -- Degree programme name
\newcommand{\ProgrammeNameValue}{Multilingual Technologies}

% -- Supervisors
\newcommand{\SupervisorValue}{\textgreater degree first name family name \textless}
\newcommand{\CoSupervisorValue}{}

% -- Commands to update the values
\newcommand{\Title}[1]{\renewcommand{\TitleValue}{#1}}
\newcommand{\Subtitle}[1]{\renewcommand{\SubtitleValue}{#1}}
\newcommand{\Name}[1]{\renewcommand{\NameValue}{#1}}
\newcommand{\Degree}[1]{\renewcommand{\DegreeValue}{#1} }
\newcommand{\Year}[1]{\renewcommand{\YearValue}{#1}}
\newcommand{\ProgrammeName}[1]{\renewcommand{\ProgrammeNameValue}{#1}}
\newcommand{\Supervisor}[1]{\renewcommand{\SupervisorValue}{#1}}
\newcommand{\CoSupervisor}[1]{\renewcommand{\CoSupervisorValue}{#1}}
\newcommand{\CoSupValue}{1}

% -- Format and setup for titlepage
\newcommand{\Titlepage}{
    \begin{titlepage}
    \fontfamily{phv}\selectfont % Arial look-alike
    \begin{tikzpicture}[remember picture,overlay]
    % logo
    \node[anchor=north west,inner sep=0pt] at (current page.north west)
    {\includegraphics[width=21.19cm,height=29.98cm]{figures/HCW_Titlepage_Background.jpg}};
     
    % thesis type, title, subtitle
    \node[anchor=north west,inner sep=0pt,text width=13cm,align=left]
        at ($(current page.north west)+(4.25cm,-5.98cm)$) {
        \fontsize{20}{26.6}\selectfont Master's Thesis\par
        \vspace{32pt}
        \fontsize{14}{16.1}\selectfont\bfseries \TitleValue\par
        \vspace{8pt}
        \fontsize{12}{13.8}\selectfont\mdseries \SubtitleValue\par
    };
 
    % author, degree, programme, supervisors
    \node[anchor=north west,inner sep=0pt,text width=13cm,align=left]
        at ($(current page.north west)+(4.25cm,-20.07cm)$) {
        \fontsize{12}{16}\selectfont
        Submitted by \textbf{\NameValue}\par
        \vspace{15pt}
        To obtain the academic degree of\\
        \textbf{\DegreeValue}\par
        \vspace{81pt}
        \fontsize{10}{13.3}\selectfont
        Vienna, \YearValue\par
        \vspace{8pt}
        \begin{tabular}{@{}p{5.3cm}@{}p{7.7cm}@{}}
        Study program: & \ProgrammeNameValue \\[8pt]
        at: & Hochschule Campus Wien \\[8pt]
        supervised by: & \SupervisorValue
            \ifnum\CoSupValue=1 \newline \CoSupervisorValue\fi
        \end{tabular}
    };
 
    \end{tikzpicture}
    \end{titlepage}
}



% Important update for glossaries, before document 
\loadglsentries{meta/acronyms.tex}
\loadglsentries{meta/glossary.tex} 
\makeglossaries %https://www.overleaf.com/learn/latex/glossaries

\begin{document}

%Start with front matter

% Add your details to this tex file
% --CHANGE THE TITLEPAGE HERE

%Please use the following commands and fields to fill in the title page with your information. 
% -- Title Page
% A title that exceeds one line should be separated with \\ within the \Title{More than one \\ line} command
\Title{I wish I had a  Master's Degree} 
\Subtitle{How great that would be}
% For a longer name use // again and make sure you add all of your titles!
\Name{First Name Last Name  Title}
\Degree{Master of Science (MSc)}
\Year{2026}
\Supervisor{My favorite}
%\CoSupervisor{My second favorite}


\Titlepage %This generates the titlepage 

\chapter*{Acknowledgements}
(This section is optional)

\chapter*{Abstract}
\chapter*{Zusammenfassung}

	
\thispagestyle{empty}
\tableofcontents{}
\listoftables 
\listoffigures



%add your chapters here
\mainmatter{}

\chapter{Introduction}
This section introduces how to use this template and which format should be used for the master's thesis in Multilingual Technologies. \textbf{Please read and follow this document carefully!}. If you write your thesis with a supervisor working at Hochschule Campus Wien University of Applied Sciences, you need to include the title page by leaving the three commands shown below as they are now in the frontmatter: 

\begin{verbatim}
\input{meta/titlepage_master_fh_campus_wien}
% --CHANGE THE TITLEPAGE HERE

%Please use the following commands and fields to fill in the title page with your information. 
% -- Title Page
% A title that exceeds one line should be separated with \\ within the \Title{More than one \\ line} command
\Title{I wish I had a  Master's Degree} 
\Subtitle{How great that would be}
% For a longer name use // again and make sure you add all of your titles!
\Name{First Name Last Name  Title}
\Degree{Master of Science (MSc)}
\Year{2026}
\Supervisor{My favorite}
%\CoSupervisor{My second favorite}


\Titlepage %This generates the titlepage 
\end{verbatim}

Replace the entries in the file \verb|titlepage_details.tex| with your own data to complete the title page. The page \textit{Declaration of Authorship} is also only needed for Thesis written at HCW.

If you write your thesis with a supervisor employed at the University of Vienna, the title page will be auto-generated in the online submission system. Please comment out the three lines shown above in the frontmatter of this file to remove the title page. 

At the University of Vienna master's thesis completed in English \textbf{have to contain a German summary}. So please make sure to complete the English abstract and German Zusammenfassung at the very beginning of this file. Both need to be submitted as part of your thesis. 


\chapter{Formatting Guidelines}
All chapters and subchapters need to be numbered continuously, which you can achieve by using \textit{chapter} for the main chapter and \textit{section} for subchapters. 

\section{References and Bibliography}
The bibliography is automatically provided for all publications and works you cite in the \verb|main.tex| file in  LaTex. You can add new entries in the \verb|bibliography.bib| file. To have a complete bibliography, the bib entries  need to contain \textbf{all required information}. So please make sure that your bib entries in the bib file are complete.  To cite papers, please use the \citep{jurafskyspeech} or \citet{jurafskyspeech}command. If you wish or have to include page numbers, add square brackets, e.g. \citep[][123]{jurafskyspeech}. 

To reference sections, figures, tables, equations, images, etc., you need to introduce so-called labels to the respective elements. In general, labels are introduced right after beginning or introducing the element. So, for instance, for sections, the label is introduced right after the \verb|\section{}| command. The label can then be directly addressed and referenced by using the \verb|\ref| command, e.g. Chapter \ref{sec:format_text} or Fig. \ref{fig:figure example}.


\section{Text}
\label{sec:format_text}
Each paragraph apart from those right after a heading should be indented. Fortunately, this is done automatically in LaTex if you introduce a blank line after the paragraph. 

This is a new paragraph. To introduce a literal quote, please use quotation marks of the style ``This is a literal quote'' \citep{lamport94} and make sure to always add number of the page on which you found the quote in the original work, e.g. \citep[][20]{lamport94}. If the quote extends over more than three lines, please use the explicit \textit{quote} environment instead of quotation marks: 

\begin{quote}
This is of a longer quote that extends over more than three lines and also needs to be equipped with a direct reference as well as page number. \citep[][8]{Goodfellow-et-al-2016}
\end{quote} 

\noindent Keep in mind that all literal quotes need to correspond 100\% to the original, character for character, including the type of text encoding. If you wish to continue the previous paragraph after a quote, table or figure, please use the \verb|\noindent| command as shown at the beginning of this paragraph.

For bullet lists you can use the environment \textit{itemize}: 
\begin{itemize}
    \item This is a first list item 
    \item And this is a second 
    \begin{itemize}
        \item And lists can even be nested
    \end{itemize}
\end{itemize}

To introduce a numbered list, please use the environment \textit{enumerate}:
\begin{enumerate}
    \item This is the first numbered item 
    \item This is the second item
\end{enumerate}

You might also want to introduce equations, for which you should use the environment \textit{equation} and a label so that you can explicitly reference the equation in your text, such as Equation~\ref{eqn:einstein}.

\begin{equation}
    \label{eqn:einstein}
     E=mc^2
\end{equation}

Any websites or URLs should \textbf{never} be simply written as text. Please add them as a bibentry to be referenced in the bibliography. If, for some reason, there is still a need to reference a URL in text, then make sure to use the command \texttt{url} and add it as a footnote as in \footnote{\url{https://www.overleaf.com/learn/latex}}.

\section{Tables and Images}
The advantage of LaTex is that it allows for continuous numbering of figures and tables and facilitates the layout of figures and tables. 

In line with the institute's guidelines, each figure and table needs to have a caption. For tables the captions need to be \textit{above} the table, as exemplified in Table \ref{tab:table example}. Please be aware of the fact that the label needs to be introduced after the caption to be able to refer tables in text. Each table and figure has to be referenced in text, e.g. Table~\ref{tab:table example}.

\begin{table}[htb]
    \centering
    \caption{Some description of the table}
	\label{tab:table example}
    \begin{tabular}{|l|c|}
    \hline
    \textbf{Command} & \textbf{Output}\\
    \hline
    \verb|{\"a}| & {\"a} \\
    \verb|{\^e}| & {\^e} \\
    \verb|{\`i}| & {\`i} \\ 
    \verb|{\.I}| & {\.I} \\ 
    \verb|{\o}| & {\o} \\
    \verb|{\'u}| & {\'u}  \\ 
    \verb|{\aa}| & {\aa}  \\\hline
    \end{tabular}
\end{table}

For any kind of figure, that is, diagrams, pictures, graphical representation, the caption needs to be introduced \textit{below} the figure. For example, see Fig. \ref{fig:figure example}

\begin{figure}[ht]
\centering
\includegraphics[scale=0.5]{figures/logo_uni_wien.png}
\caption{Some description of the figure}
\label{fig:figure example}
\end{figure}

\section{Code}
To include code in your text, you should use the environment \textit{verb} as in \verb|code|.

For a longer block of code, please use the verbatim environment: 

\begin{verbatim}
Text enclosed inside \texttt{verbatim}
environment is printed directly 
and all \LaTeX{} commands are ignored.
\end{verbatim}

Alternatively, you can also use the \textit{lstlisting} environment to encode longer blocks of code, which you can use as is or format and color-code as you see fit. 

\begin{lstlisting}[language=Python,caption={Python example}]
from transformers import AutoTokenizer

tokenizer = AutoTokenizer.from_pretrained("google-bert/bert-base-cased")

encoded_input = tokenizer("Do not meddle in the affairs of wizards, for they are subtle and quick to anger.")
print(encoded_input)

\end{lstlisting}

\section{Algorithms}
When representing algorithms in LaTex it is common to represent them as pseudcode\footnote{See \url{https://www.overleaf.com/learn/latex/Algorithms} for more details.}. Please also make sure to reference each included algorithm explicitly in the text, e.g. Algorithm~\ref{alg:cap}.

\begin{algorithm}
\caption{An algorithm with caption}\label{alg:cap}
\begin{algorithmic}
\State $i \gets 10$
\If{$i\geq 5$} 
    \State $i \gets i-1$
\Else
    \If{$i\leq 3$}
        \State $i \gets i+2$
    \EndIf
\EndIf 
\end{algorithmic}
\end{algorithm}

\section{Acronyms and Glossary}

If you want to use acronyms or a glossary check the page here: \url{https://www.overleaf.com/learn/latex/glossaries}

The \Gls{latex} typesetting markup language is specially suitable 
for documents that include \gls{maths}. are 
rendered properly an easily once one gets used to the commands.

Given a set of numbers, there are elementary methods to compute 
its \acrlong{gcd}, which is abbreviated \acrshort{gcd}. This 
process is similar to that used for the \acrfull{lcm}.


\bibliography{bibliography}
\bibliographystyle{acl_natbib}
%\printbibliography


\printglossary[type=\acronymtype]
\printglossary

\cleardoublepage{}
\pagebreak
	
%add your end matter here
\appendix{}

\chapter{Appendix}
	
\end{document}